\documentclass[11pt]{article}
\usepackage[utf8]{inputenc}
\usepackage{amsmath,amssymb}
\usepackage[margin=1in]{geometry}
\usepackage{hyperref}
\emergencystretch=3em

\title{The real weighted box integral: mass, recursion, permutations}
\author{Claude, with Daniel Murfet}
\date{2026-09-07}

\begin{document}
\maketitle
\sloppy

\begin{abstract}
Third step of the mixed-ratio general-$d$ monomial programme (Astra~\#15, route~A). The real
weighted box integral
$W_\ell(N)=\int_{(0,1]^d}\prod_it_i^{\ell_i-1}e^{-\beta N\prod_it_i}\,dt$ is defined as the
\texttt{toReal} of the $[0,\infty]$-valued \texttt{weightedBoxIntegral}, and the four facts the
mixed-ratio induction needs are proved: the mass bound
$\int_{(0,1]^d}\prod_it_i^{w_i}\,dt=\prod_i1/(w_i+1)$ (hence $0\le W_\ell(N)\le\prod_i1/\ell_i$ for
$\beta N\ge0$); measurability of $N\mapsto W_\ell(N)$; the real recursion
$W_\ell(N)=\int_0^1t^{\ell_0-1}W_{\ell'}(Nt)\,dt$ with $\ell'$ the tail; and permutation invariance
of the weighted box integral. Zero \texttt{sorry}/\texttt{axiom}.
\end{abstract}

\section{The things formalised}
{\small
\begin{verbatim}
theorem weightedBoxIntegral_one : ∀ d w, (∀ i, -1 < w i) →
    weightedBoxIntegral d w (fun _ => 1) = ofReal (∏ i, 1 / (w i + 1))
theorem weightedBoxIntegral_expKernel_lt_top (d w) (hw) (β N) (hβN : 0 ≤ β * N) :
    weightedBoxIntegral d w (expKernel β N) < ⊤
def mixedBoxReal (d) (ℓ : Fin d → ℝ) (β N : ℝ) : ℝ :=
  (weightedBoxIntegral d (fun i => ℓ i - 1) (expKernel β N)).toReal
theorem mixedBoxReal_le (d ℓ) (hℓ : ∀ i, 0 < ℓ i) (β N) (hβN : 0 ≤ β * N) :
    mixedBoxReal d ℓ β N ≤ ∏ i, 1 / ℓ i
theorem measurable_mixedBoxReal (d ℓ β) : Measurable fun N => mixedBoxReal d ℓ β N
theorem mixedBoxReal_succ (d) (ℓ : Fin (d+1) → ℝ) (hℓ) (β N) (hβ : 0 < β) (hN : 0 < N) :
    mixedBoxReal (d + 1) ℓ β N
      = ∫ t in Ioc 0 1, t ^ (ℓ 0 - 1) * mixedBoxReal d (Fin.tail ℓ) β (N * t)
theorem weightedBoxIntegral_comp_perm (d w g) (σ : Equiv.Perm (Fin d)) :
    weightedBoxIntegral d (w ∘ σ) g = weightedBoxIntegral d w g
\end{verbatim}
}

\section{Mathematics}
The mass bound is the constant-kernel case of the weighted recursion,
$\int_0^1a^{w_0}\,da=1/(w_0+1)$ at each level; since $e^{-\beta Nz}\le1$ on the box for
$\beta N\ge0$, the exponential integral is finite and bounded by the mass. The real recursion is the
$[0,\infty]$ recursion of unit~174 after \texttt{integral\_toReal}, which needs the parametric
integral $M\mapsto W(M)$ to be measurable (a Tonelli-type statement,
\texttt{Measurable.lintegral\_prod\_right'}) and finite almost everywhere on $(0,1]$. Permutation
invariance transports the integral along \texttt{MeasurableEquiv.piCongrLeft}, which preserves the
product of restricted Lebesgue measures, and reindexes the two products with
\texttt{Equiv.prod\_comp}.

\section{Paper-facing meaning}
$W_\ell$ is the substituted normal moment integral with exponents $\ell_i=(h_i+1)/(2k_i)$; the
recursion peels one normal coordinate, and permutation invariance lets the induction of the next
unit always peel a coordinate that does \emph{not} realise the minimal ratio, which is where the
dominated-coordinate transfer lemma applies.

\section{Lean notes}
\begin{itemize}
\item \texttt{positivity} cannot prove $0\le1/(w+1)$ from a hypothesis $-1<w$; use
\texttt{div\_nonneg zero\_le\_one (by linarith)}.
\item \texttt{lintegral\_mul\_const} with a composed measurability proof leaves a
\texttt{(ofReal ∘ f) a} pattern that does not match the goal; \texttt{lintegral\_mul\_const'} with
\texttt{ENNReal.ofReal\_ne\_top} avoids measurability altogether.
\item \texttt{mul\_le\_mul\_left'} is gone on this pin; \texttt{gcongr} handles
$a\cdot b\le a\cdot c$ in \texttt{ℝ≥0∞}.
\item \texttt{Equiv.piCongrLeft\_apply\_apply} takes the family and the equivalence explicitly:
\texttt{Equiv.piCongrLeft\_apply\_apply (fun \_ => ℝ) σ y b}.
\item \texttt{rw [restrict\_unitBox]} rewrites both occurrences at once; a second \texttt{rw} fails.
\item Rewrite the measure-preserving transport on the right-hand side with \texttt{conv\_rhs} so
that the pointwise goal has the composed integrand on the side where the reindexing lemma applies.
\end{itemize}

\end{document}
